% !TeX root = ../../main.tex
\section{分层执行的数据一致性与资源约束}
\subsection{离线准备的数据一致性}
两阶段系统同时维护多个矩阵视图，其一致性首先依赖行标识。设原始行号为 $i$，重排后行号为 $\pi(i)$。Q-core 输出的编号通常对应重排后存储位置，Reader 必须按同一编号体系查找 FP32 副本，最终结果再通过逆映射恢复原始条目。若两个副本独立重排，即使粗筛得分正确，重算也可能读取错误对象。

因此，预处理应建立唯一的重排向量，并使量化、ReSparse、格式编码和原始行映射共享这一向量。对每个候选，验证可比较其原始行号、非零数量及列索引序列，而不是仅验证最终输出数量。原始矩阵还应保持只读，使离线剪枝不会意外污染重算数据。

\begin{table}[htbp]\centering\caption{AccelES 各数据视图的一致性要求}
\begin{tabularx}{\textwidth}{p{2.6cm}XX}\toprule
数据视图&应保留的信息&不一致时的影响\\\midrule
原始 FP32 矩阵&原始数值、列号、行标识&重算失去参考语义\\
ReSparse 矩阵&保留位置与原行对应&候选评价与硬件输入不同\\
INT6 矩阵&量化尺度和零元素处理&软硬件得分不一致\\
Ultra-CSR 数据包&行边界和跨包延续&行得分被拆分或错误合并\\
Random-CSR 元数据&首包、偏移、长度&读取越界或混入相邻行\\
逆置换表&原条目与存储行映射&返回 ID 无法对应真实候选\\\bottomrule
\end{tabularx}\end{table}

\subsection{定点累加位宽与溢出控制}
INT6 输入减少乘法器位宽，但累加器仍须容纳整行和。两个有符号6 位整数的乘积绝对值最大可达1024，若一行有 $n_i$ 个参与乘法的元素，则最坏情况绝对和不超过 $1024n_i$。采用 $b_a$ 位有符号累加器时，应满足
\begin{equation}
 2^{b_a-1}-1\geq1024n_i.
\end{equation}
例如仅作上界分析，$n_i\leq1024$ 时使用不少于22 位的有符号累加器可以覆盖这一绝对值上界。该上界给出了累加位宽的正确性约束，具体硬件配置还需结合资源与时序确定。

若实际输入只使用 $[-31,31]$，乘积界可收紧为961，但第三章量化允许 $-32$，应采用与真实编码一致的界。跨包部分和与最终行和必须使用相容位宽，不能只为单包30 个乘积的和分配空间。饱和与回绕也有不同错误语义，应在软硬件仿真中显式区分。

低位宽乘法通过 DSP 打包实现时，需要防止子乘积之间的进位污染。不同打包策略可能对符号扩展和保护位提出额外要求，因此“一个 DSP 支持多路乘法”是具体实现能力，而非任意有符号操作都可无条件复用。资源节省必须与逐元素正确性验证共同检查。

\subsection{跨包边界的状态维护}
连续流式数据包可能包含一行的尾部、若干完整短行以及下一行的头部。归约器应区分完整行输出与未完成行缓存。若当前包首元素继续上一包末行，先合并保存的部分和；若首元素已切换到新行，则先提交上一行，再开始新累加。

设包 $t$ 的末行标识为 $r_t$、未完成和为 $s_t$。包 $t+1$解码后，只有属于同一行的乘积才能与 $s_t$ 合并。一次传输中可能结束多行，候选更新单元的吞吐也需匹配多行完成率。只保证每周期读取30 个元素，却不能及时接收多个行得分时，反压仍会降低实际带宽。

\begin{table}[htbp]\centering\caption{流式边界情况及对应处理}
\begin{tabularx}{\textwidth}{p{3cm}XX}\toprule
情况&处理要求&验证重点\\\midrule
单行跨多个包&持续保存部分和&不能提前输出或丢弃中间包\\
一包内含多个短行&分别归约并提交&多行输出是否堵塞流水\\
包首恰为新行&提交前包末行&避免跨行累加\\
最后一包未填满&屏蔽填充记录&避免产生虚假得分和候选\\
量化后空行&保持行号映射&边界标记能否表达空行跳跃\\
负得分占多数&使用正确候选初值&不可用零值替代负无穷边界\\\bottomrule
\end{tabularx}\end{table}

这些情况可以通过小规模可手算输入验证，独立于大数据集性能测量。小输入应覆盖同一行跨越两个包、整包恰好结束一行以及全负得分等情况。它们的作用是检测控制逻辑，而非证明大规模候选覆盖。

\subsection{局部候选容量的概率解释}
随机重排行缓解重要候选集中在同一核的风险。若作理想化假设，真实 Top-K 均匀独立分配到 $P$ 个分区，则某个分区包含的真实重要候选数近似服从二项分布 $X\sim\mathrm{Binomial}(K,1/P)$。局部容量为 $l$ 时，该分区超容量的概率为
\begin{equation}
 \Pr(X>l)=1-\sum_{j=0}^{l}{K\choose j}(1/P)^j(1-1/P)^{K-j}.
\end{equation}
使用并合界，任意分区超容量概率不超过 $P\Pr(X>l)$。这只是随机分配下的辅助解释，重要候选之间的相关性、实际重排和近似得分排序会使独立假设不成立。

因此，不能用该概率表达式取代实验召回。它可以说明为什么增大局部容量有助于降低分区截断，并帮助选择需要验证的候选数范围。对同一数据库应固定重排种子以便复现，同时使用多个种子检查结果是否对偶然分区敏感。

\subsection{两阶段吞吐与资源分配}
若每个查询的粗筛耗时为 $T_Q$、重算耗时为 $T_F$，不重叠时在线计算成本近似为两者之和。若多个查询能够使用双缓冲形成阶段流水，稳态吞吐可能受 $\max(T_Q,T_F)$限制，但单次查询延迟仍包含其经过各阶段的时间。必须区分本研究已报告的执行方式与这种潜在流水扩展，不能把理想稳态吞吐写成实测结果。

高精度核心共享能够降低资源占用，其合理性依赖候选数据量较小。随着 $L$增加、候选行变长或查询并发提高，F-core 可能饱和。此时新增 Q-core 不会提高系统吞吐，反而扩大候选队列。资源分配宜观察各阶段忙闲比例、队列深度和反压时间，而不是只比较 DSP 利用率。

对宽度为 $W$的传输通道，若候选平均需要 $b$ 个包，每核输出 $l$个候选，则第二阶段至少涉及约 $Plb$个数据包，加上行元数据读取。该数量与第一阶段完整矩阵包数比较，可估计重算是否仍属于小规模任务。真实随机事务延迟还取决于通道并发，不能直接将包数除以峰值带宽得到准确时间。

\subsection{参数校准与泛化验证}
候选预算、保留比例和量化位宽应在校准查询上选择，在独立查询上报告质量。若使用同一批查询调参并汇总，结果可能低估新查询损失。对于数据分布差异较大的语料，应分别记录参数，而不将某个语料上的最优配置当作统一常数。

校准可先固定较宽候选预算验证 INT6，再逐步加入局部候选截断和 ReSparse，最后检查完整组合。顺序加入模块有助于发现损失来源，但每个模块的独立作用还需与全部配置做交叉验证，因为量化与剪枝可能存在相互影响。总体质量损失不能简单等于各模块单独损失之和。

\subsection{更新场景与接口边界}
数据库插入新行时，需要同步生成低比特数据和 FP32 行元数据；修改现有行则可能改变包边界和分区大小。静态构建的连续数据包不一定适合原地插入，因此更可行的扩展包括增量缓冲和周期重构。本文未报告这种动态维护系统的实测性能，属于部署时需要补充的问题。

查询接口应明确返回候选 ID、重算得分还是最终 Top-K。若只传输 ID，主机需要后续访问原始数据时，应计入完整应用成本；若传输得分，则需注明其数值精度和对应缩放。对用户而言，正确的接口规范比单一内核吞吐更能决定系统能否直接替换原有实现。
